\chapter{Serum Free Light Chains}

\section*{What Is This Test?}
The serum free light chain (FLC) assay measures the concentration of kappa ($\kappa$) and lambda ($\lambda$) immunoglobulin free light chains that circulate in serum in an unbound state --- that is, not bound to immunoglobulin heavy chains. In normal B-cell biology, plasma cells produce intact immunoglobulins, and a slight excess of free light chains (approximately 40\% excess) is synthesized relative to heavy chains. These free light chains are secreted into the blood. Due to their small size (kappa monomers ~22.5 kDa; lambda dimers ~45 kDa), free light chains are rapidly cleared from the circulation by the kidneys: they are freely filtered through the glomerulus, reabsorbed by proximal tubular cells via the cubilin-megalin receptor complex, and catabolized within the tubular cells. The serum half-life of free light chains is short: 2--4 hours for kappa, 3--6 hours for lambda (compared to 20--23 days for intact IgG). The FLC assay is essential for the diagnosis, classification, and monitoring of plasma cell dyscrasias. It detects monoclonal free light chain production that may be missed by serum protein electrophoresis (SPEP), particularly in light chain multiple myeloma (15--20\% of myelomas), AL amyloidosis, and light chain deposition disease. The kappa/lambda ($\kappa/\lambda$) FLC ratio is a sensitive indicator of clonality. The assay has largely replaced the 24-hour urine protein electrophoresis (UPEP) for initial screening of plasma cell disorders, though UPEP remains important for quantifying significant proteinuria and detecting light chain cast nephropathy \cite{katzmann2006elimination}.

\section*{Alternative Names}
FLC; Serum Free Light Chain Assay; Freelite; Kappa Free Light Chains; Lambda Free Light Chains; Kappa/Lambda Ratio; $\kappa/\lambda$ Ratio; FLC Ratio; Light Chain Assay

\section*{Principle}
The serum free light chain assay (Freelite, The Binding Site/Thermo Fisher) uses a latex-enhanced immunoturbidimetric or immunonephelometric method. The assay employs polyclonal sheep antibodies that specifically recognize epitopes on free kappa or free lambda light chains that are hidden (masked) when the light chain is bound to the immunoglobulin heavy chain in intact immunoglobulins. These epitopes are exposed on the constant domain of the light chain only when the light chain is free (unbound). Patient serum is incubated with latex particles coated with anti-free kappa or anti-free lambda specific antibodies. The resulting agglutination increases turbidity, which is measured photometrically. The concentration of free light chains is proportional to the turbidity and is calculated from a calibration curve. Results are reported in mg/L for free kappa, free lambda, and the calculated $\kappa/\lambda$ ratio. The monoclonal immunoglobulins can interfere with the assay if present at very high concentrations through non-specific binding or matrix effects. High-dose hook effect (false-negative due to antigen excess) must be considered when free light chain levels are unexpectedly normal in a patient with suspected light chain disease; the laboratory should perform serial dilutions of the sample to unmask the hook effect \cite{dispenzieri2009international}. A newer alternate methodology uses mass spectrometry (MALDI-TOF) to identify and quantify monoclonal light chains based on their unique molecular mass.

\section*{History}
The traditional method for detecting monoclonal free light chains was the Bence Jones protein test, first described by Henry Bence Jones in 1848. The modern serum free light chain (Freelite) assay was developed by Arthur Bradwell and colleagues at the University of Birmingham, UK, and introduced into clinical practice in 2001. The assay transformed the workup of monoclonal gammopathies by allowing detection of serum free light chains with far greater sensitivity than SPEP and UPEP. The International Myeloma Working Group (IMWG) incorporated the FLC assay into the diagnostic criteria for multiple myeloma in 2006 and further refined its use in the 2014 updated diagnostic criteria \cite{rajkumar2014international}. The FLC assay is now a mandatory component of the initial workup of all monoclonal gammopathies and is recommended alongside SPEP, immunofixation, and 24-hour UPEP.

\section*{Why This Test Is Ordered}
\begin{itemize}
  \item Screening for monoclonal gammopathies, in conjunction with serum protein electrophoresis (SPEP) and immunofixation (IFE). The FLC assay detects light chain production missed by SPEP and is essential for identifying light chain multiple myeloma, AL amyloidosis, and light chain deposition disease
  \item Diagnosis of light chain multiple myeloma (Bence Jones myeloma), which accounts for 15--20\% of all myelomas. In these patients, intact immunoglobulin is not produced; the M-protein may not be visible on SPEP because only free light chains are secreted. The FLC assay (with an abnormal $\kappa/\lambda$ ratio) replaces 24-hour UPEP for initial diagnosis
  \item Diagnosis and monitoring of AL amyloidosis (immunoglobulin light chain amyloidosis). The FLC assay is the most sensitive test for detecting the clonal plasma cell dyscrasia underlying AL amyloidosis (abnormal $\kappa/\lambda$ ratio is present in >90\% of patients) \cite{palladini2018serum}. Quantitative FLC monitoring is used to assess hematologic response to chemotherapy (the difference between involved and uninvolved free light chains [dFLC] is the key parameter)
  \item Prognostication in monoclonal gammopathy of undetermined significance (MGUS) --- an abnormal FLC ratio is one of three risk factors for progression of MGUS to multiple myeloma (the other two are non-IgG isotype and M-protein >1.5 g/dL). MGUS with all three risk factors has a 58\% risk of progression at 20 years
  \item Monitoring of treatment response in multiple myeloma and AL amyloidosis. The IMWG uniform response criteria use the FLC assay to define stringent complete response (sCR), complete response (CR), and very good partial response (VGPR). In AL amyloidosis, a dFLC <40 mg/L defines a complete hematologic response
  \item Diagnosis of monoclonal gammopathy of renal significance (MGRS) --- a spectrum of kidney diseases caused by nephrotoxic monoclonal immunoglobulins or free light chains in patients who do not meet criteria for overt multiple myeloma
  \item Baseline evaluation and monitoring of all patients with newly diagnosed multiple myeloma before, during, and after therapy
\end{itemize}

\section*{Primary vs Secondary Test}
This is a primary diagnostic test for plasma cell dyscrasias. The IMWG recommends that serum FLC assay be performed together with SPEP and immunofixation as the initial screening panel for all suspected monoclonal gammopathies. In many circumstances, the FLC assay has replaced the 24-hour UPEP for initial evaluation.

\section*{How This Test Is Performed}
For most patients, a health care professional will take a blood sample from a vein in your arm using a small needle. After the needle is inserted, a small amount of blood will be collected into a test tube or vial. You may feel a little sting when the needle goes in or out. This usually takes less than five minutes.

\section*{What Preparation Is Needed}
No special preparation is required. Fasting is not necessary. Serial FLC measurements should be performed using the same assay and laboratory to minimize variability. Patients should inform their healthcare provider if they have impaired renal function, as renal impairment affects free light chain clearance and changes the FLC ratio reference range.

\section*{Contraindications and Precautions}
No absolute contraindications. The most critical limitation is that the FLC $\kappa/\lambda$ ratio reference range (normally 0.26--1.65) is valid only for patients with NORMAL renal function. In renal impairment, the $\kappa/\lambda$ ratio reference range expands (``renal reference range'' typically 0.37--3.1 for patients with CKD stage 3 or greater) because lambda free light chains (normally dimers, ~45 kDa) are cleared more slowly than kappa free light chains (normally monomers, ~22.5 kDa) when GFR is reduced. This results in a higher $\kappa/\lambda$ ratio in renal impairment that is POLYCLONAL (both kappa and lambda increase), not monoclonal. Failure to use the renal reference range results in false-positive diagnosis of a kappa-restricted plasma cell dyscrasia in patients with CKD. Polyclonal increases in free light chains occur in: chronic kidney disease (both chains increase, ratio normal when using renal range), systemic inflammation and autoimmune diseases (polyclonal B-cell activation), chronic infections (HIV, HCV, HBV, EBV), and lymphoproliferative disorders with polyclonal B-cell activation. Very high concentrations of intact monoclonal immunoglobulin (>5 g/dL) can interfere with the FLC measurement through non-specific binding and produce falsely abnormal results. High-dose hook effect (antigen excess causing falsely low FLC levels) can occur when the free light chain concentration exceeds the assay's linear range (typically >10,000 mg/L). Specimen should be collected in a serum separator tube (SST, gold-top) or red-top tube. Grossly hemolyzed or lipemic specimens may interfere with nephelometry.

\section*{Risks}
Minimal risk. Slight pain or bruising at the venipuncture site. Rare: hematoma, infection, vasovagal syncope.

\section*{What Happens After the Test}
Results are typically available within 1 to 3 days. An abnormal $\kappa/\lambda$ ratio indicates a clonal plasma cell disorder. The direction of the abnormality (ratio >1.65 indicates kappa restriction; ratio <0.26 indicates lambda restriction) guides the interpretation. An abnormal ratio must be interpreted together with the absolute values of free kappa and free lambda, SPEP, immunofixation, and clinical findings. A very high free light chain concentration (involved chain >500--1,000 mg/L) with a markedly abnormal ratio in a patient with acute kidney injury suggests light chain cast nephropathy and constitutes a hematologic emergency requiring urgent initiation of therapy.

\section*{Understanding Results}

\subsection*{Below Normal (Ratio <0.26 --- Lambda Restriction)}
\begin{itemize}
  \item Lambda free light chain multiple myeloma: Lambda-predominant monoclonal free light chain production with suppressed kappa levels, producing a $\kappa/\lambda$ ratio <0.26 (often <0.1)
  \item AL amyloidosis (lambda type): Lambda is the more common light chain in AL amyloidosis (approximately 70\% lambda, 30\% kappa). The ratio is typically <0.26. The difference between involved and uninvolved free light chains (dFLC) is the key monitoring parameter
  \item Lambda light chain MGUS: An abnormal $\kappa/\lambda$ ratio <0.26 with a low level of lambda free light chain, no end-organ damage, and no clinical symptoms. Requires monitoring, not treatment
  \item Lambda monoclonal gammopathy of renal significance (MGRS): Nephrotoxic lambda free light chains cause kidney injury (usually focal segmental glomerulosclerosis or membranoproliferative glomerulonephritis pattern) without meeting criteria for myeloma
  \item Immunoglobulin suppression (immunoparesis): The uninvolved free light chain (kappa, in this case) is suppressed, reflecting the disrupted normal plasma cell function in the bone marrow due to the clonal plasma cell population
\end{itemize}

\subsection*{Above Normal (Ratio >1.65 --- Kappa Restriction)}
\begin{itemize}
  \item Kappa free light chain multiple myeloma: Kappa-predominant monoclonal free light chain production (kappa accounts for approximately 65\% of light chain myelomas). Ratio >1.65, often >10--100
  \item AL amyloidosis (kappa type): Less common than lambda AL amyloidosis. $\kappa/\lambda$ ratio >1.65
  \item Kappa light chain MGUS: Abnormal $\kappa/\lambda$ ratio >1.65 with a low level of kappa free light chain, no end-organ damage
  \item Kappa monoclonal gammopathy of renal significance (MGRS)
  \item Polyclonal increase in the $\kappa/\lambda$ ratio due to RENAL IMPAIRMENT: Both kappa and lambda increase, but the $\kappa/\lambda$ ratio may be mildly elevated (up to ~3.1) due to slower clearance of lambda dimers. Critical to use the renal reference range (0.37--3.1) rather than the normal reference range (0.26--1.65) in patients with CKD (eGFR <60 mL/min). A $\kappa/\lambda$ ratio >3.1 in a patient with CKD is more likely to indicate a true kappa-restricted monoclonal process
  \item Kappa free light chain >500 mg/L with $\kappa/\lambda$ ratio >100: Very high risk for light chain cast nephropathy and acute kidney injury. This is a hematologic emergency. Requires urgent initiation of anti-plasma cell therapy (cyclophosphamide-bortezomib-dexamethasone [CyBorD] or daratumumab-based regimen) to rapidly reduce free light chain production and prevent irreversible renal damage
  \item Intact immunoglobulin multiple myeloma (IgG kappa, IgA kappa) with concomitant free light chain production: Both intact monoclonal immunoglobulin AND excess monoclonal free kappa light chain are produced. The SPEP shows an M-spike, and the FLC assay shows an abnormal ratio
\end{itemize}

\section*{Normal Results}
\begin{itemize}
  \item \textbf{Free Kappa:} 3.3--19.4 mg/L
  \item \textbf{Free Lambda:} 5.7--26.3 mg/L
  \item \textbf{Kappa/Lambda Ratio (normal renal function):} 0.26--1.65
  \item \textbf{Kappa/Lambda Ratio (renal reference range, CKD stage 3 or greater):} 0.37--3.1
  \item \textbf{Key diagnostic thresholds (IMWG):}
  \begin{itemize}
    \item An abnormal FLC ratio (ratio <0.26 or >1.65 in setting of normal renal function) is evidence of clonality
    \item An involved/uninvolved FLC ratio $\geq$100 (with involved FLC $\geq$100 mg/L) is a myeloma-defining event (SLiM CRAB criterion)
    \item The difference between involved and uninvolved free light chains (dFLC) should be used for monitoring in AL amyloidosis
  \end{itemize}
  \item The FLC ratio is more sensitive than immunofixation for detecting residual clonal disease, but it must be interpreted in the context of the absolute FLC values. A ratio that is abnormal due to profoundly low uninvolved chain (immunoparesis) rather than an elevated involved chain has different clinical significance than a ratio driven by an elevated involved chain
  \item Serial FLC measurements should use the same assay and ideally the same laboratory, as inter-assay variability can be significant
\end{itemize}

\section*{Follow-up Tests}
\begin{itemize}
  \item \textbf{Serum protein electrophoresis (SPEP) and serum immunofixation electrophoresis (IFE):} Should be ordered simultaneously with the FLC assay as part of the initial screening panel for monoclonal gammopathies (FLC + SPEP + IFE). SPEP detects and quantifies intact monoclonal immunoglobulins. IFE identifies the immunoglobulin heavy chain class and light chain type. The combination of all three tests maximizes sensitivity for detecting plasma cell disorders
  \item \textbf{24-hour urine protein electrophoresis (UPEP) and urine immunofixation:} Quantifies urinary excretion of monoclonal free light chains (Bence Jones protein). Still important for quantifying significant light chain proteinuria, particularly in AL amyloidosis with nephrotic syndrome, and for diagnosing light chain cast nephropathy in the setting of acute kidney injury \cite{dejoie2016serum}
  \item \textbf{Bone marrow biopsy and aspiration:} Essential for diagnosis. Quantifies the percentage of clonal plasma cells (diagnostic threshold $\geq$10\% for symptomatic MM, $\geq$60\% as a myeloma-defining event). Flow cytometry and cytogenetics/FISH are performed on the bone marrow aspirate
  \item \textbf{Cytogenetics and FISH (fluorescence in situ hybridization) on bone marrow:} For detection of high-risk cytogenetic abnormalities: del(17p), t(4;14), t(14;16), t(14;20), and 1q amplification. These are essential for R-ISS staging and guiding treatment decisions
  \item \textbf{Skeletal survey (whole-body low-dose CT, FDG-PET/CT, or whole-body MRI):} For detection of lytic bone lesions in multiple myeloma. MRI is more sensitive than CT for diffuse bone marrow infiltration
  \item \textbf{Serum creatinine, eGFR, and serum calcium:} To assess for the CRAB criteria (hyperCalcemia, Renal insufficiency, Anemia, Bone lesions) that define symptomatic myeloma requiring treatment. Renal impairment from light chain cast nephropathy is a medical emergency
  \item \textbf{Serum albumin, LDH, and beta-2 microglobulin:} Components of the R-ISS staging system for multiple myeloma. Beta-2 microglobulin reflects tumor burden and renal function
  \item \textbf{Abdominal fat pad biopsy, rectal biopsy, or organ biopsy (kidney, heart, liver):} For diagnosis of AL amyloidosis. Congo red staining with apple-green birefringence under polarized light is diagnostic. Mass spectrometry-based proteomic analysis determines the amyloid type (AL, ATTR, AA, etc.)
  \item \textbf{NT-proBNP or BNP, troponin T or I, and echocardiogram:} For assessment of cardiac involvement in AL amyloidosis. NT-proBNP and troponin T form the basis of the Mayo Clinic AL amyloidosis cardiac staging system
\end{itemize}

\section*{References}
\bibliographystyle{unsrtnat}
\bibliography{references}

\clearpage
\section*{Regulatory Status and Authorization Date}
Regulatory status applies to the specific assay, instrument, manufacturer, and intended use rather than to the test name alone. No single approval date can be assigned to this test category. For a commercial assay, verify the current product in the relevant regulator's database: the U.S. Food and Drug Administration (FDA) clearance, approval, or authorization; India's Central Drugs Standard Control Organisation (CDSCO) licence or permission; the European Union In Vitro Diagnostic Regulation (EU IVDR) CE certificate issued through the applicable conformity-assessment route; or the United Kingdom Medicines and Healthcare products Regulatory Agency (MHRA) registration and UKCA/CE status. Laboratory-developed tests may not have an individual FDA, CDSCO, EU IVDR, or MHRA approval date.
\clearpage
